\documentclass[preamble=combinatory-preamble.tex]{markdown}
title: A combinatory formulation of PTS-rk
date: October 22, 2026

:::Abstract
PTS-rk admits a complete Hilbert presentation by finitely many numerically
indexed axiom schemes, by the supersorted-system theorem of Bunder and
Dekkers. A finite basis of fixed closed typed combinators is nevertheless
impossible: application cannot increase the rank of its function's type.
We give rank-typed BCKW arithmetic, an explicit Hilbert catalogue and a
first-order algebraic presentation using quotient
inductive-inductive-recursive syntax. Checked soundness and extraction
are distinguished from the remaining formal completeness and adequacy
proofs. All expansions retain the original annotations and use beta
conversion only. Elementary bounds control normalization and every
intermediate term; fixing the input's universe labels gives finite syntax
carriers and a total, typed Agda apply-and-normalize function.
Independent certificate checking gives contextual soundness in graded finite
carriers and external exclusion of a closed proof of Void.
Finite Church encodings give explicit rank-bounded code-to-code
reduction programs, separated from the proposed internal consistency proof.
A ground-query Horn reduction gives a precise, weaker sense of
reduction to a locally finitely satisfiable theory.
:::

# The calculus and its logical reading {#sec-calculus}

Terms of PTS-rk are variables, universes $U_i$ ($i\in\mathbb N$), products
$\Pi x:A.B$, annotated abstractions $\lambda x:A.t$, and applications $tu$.
Contexts are lists of typed variables. Its distinctive rules are

$$
\Gamma\vdash U_i:U_0,\qquad
\frac{\Gamma\vdash A:U_a\quad\Gamma,x:A\vdash B:U_b}
 {\Gamma\vdash\Pi x:A.B:U_{r(a,b)}},
\qquad r(a,b)=\max(a+1,b).
$$

Typing and conversion are those of the original calculus: annotations are
retained, with neither cumulativity nor eta conversion. In particular,
$U_i:U_0$ is not the usual ascending universe rule $U_i:U_{i+1}$.

A logical judgment $\Gamma\vdash t:A$ reads: under assumptions $\Gamma$,
$t$ is a proof of the proposition represented by $A$. Types may also
describe data; there is no separate primitive universe of propositions.
A product $\Pi x:A.B(x)$ is universal quantification over inputs $x:A$;
when the input and output are propositions and $B$ does not depend on $x$,
it is implication $A\to B$. Quantification over $A:U_a$ allows types and
propositions themselves as parameters.

Application is dependent elimination:

$$
\frac{\Gamma\vdash f:\Pi x:A.B(x)\quad\Gamma\vdash u:A}
 {\Gamma\vdash fu:B(u)}.
$$

For $A\to B$, this is modus ponens, with $fu$ recording the proofs used.
Natural deduction discharges assumptions by abstraction; a Hilbert system
encodes discharge through closed axioms and application. A derivation
certifies the judgment, while its subject term is the object-level proof.
Beta reduction computes within that proof without changing its type.

# Rank-typed BCKW arithmetic {#sec-bckw}

The standard basis separates composition $B$, exchange $C$, weakening $K$
and duplication $W$; we retain identity $I$ as an additional atom.
Their types and computation rules are

$$
\begin{array}{c|c|l}
 &\text{type}&\text{beta rule}\\ \hline
B&(Y\to Z)\to(X\to Y)\to X\to Z&Bfgx\longrightarrow_\beta^*f(gx)\\
C&(X\to Y\to Z)\to Y\to X\to Z&Cfxy\longrightarrow_\beta^*fyx\\
K&X\to Y\to X&Kxy\longrightarrow_\beta^*x\\
W&(X\to X\to Y)\to X\to Y&Wfx\longrightarrow_\beta^*fxx\\
I&X\to X&Ix\longrightarrow_\beta^*x .
\end{array}
$$

Each atom expands to the corresponding annotated lambda term. Type
instantiations are implicit in these formulas, but retained in the checked
syntax; arrow ranks are computed by $r(a,b)=\max(a+1,b)$, not by an
untyped interpretation. The $B,C,K,I$ fragment is affine in its value
arguments: their computation rules do not duplicate them; $W$ does [7].
This does not require type parameters to occur affinely in annotations.

Fix $X=U_i$, put $E=X\to X$ and $\mathsf{Ch}=E\to E$, and define

$$
H=C(BB(CB))B,\qquad
\overline0=KI,\qquad
\overline{n+1}=W(H\overline n).
$$

Here $Hnfgx\longrightarrow_\beta^*f(ngx)$. Thus the first codes are

$$
\overline0=KI,\qquad
\overline1=W(H(KI)),\qquad
\overline2=W(H(W(H(KI)))).
$$

After expanding $H$ without sharing, $\overline n$ contains exactly
$n$ occurrences of $W$ and **$7n+2$ combinator atoms in total**:
$H$ has six atoms, and each successor layer adds one $W$.
The count treats $I$ as one atom and does not count implicit type indices.
For this family those indices have fixed shapes independent of $n$.

Successor, addition and multiplication are the following typed codes,
with independent type instantiations at each occurrence of an atom:

$$
\begin{aligned}
\mathsf{succ}&=BWH
  &&:\mathsf{Ch}\to\mathsf{Ch},\\
\mathsf{add}&=B\bigl(B(BW)(BBC)\bigr)(BB)
  &&:\mathsf{Ch}\to\mathsf{Ch}\to\mathsf{Ch},\\
\mathsf{mul}&=B
  &&:\mathsf{Ch}\to\mathsf{Ch}\to\mathsf{Ch}.
\end{aligned}
$$

Their atom counts are $8,9,1$, respectively. Applied to Church arguments,
their expansions satisfy

$$
\begin{aligned}
\mathsf{succ}\;\overline n\,f\,x&\longrightarrow_\beta^*f^{n+1}x,\\
\mathsf{add}\;\overline m\;\overline n\,f\,x
  &\longrightarrow_\beta^*f^{m+n}x,\\
\mathsf{mul}\;\overline m\;\overline n\,f\,x
  &\longrightarrow_\beta^*f^{mn}x.
\end{aligned}
$$

:::Theorem[Rank-typed BCKW arithmetic]
For simple-type instances over $U_i$, every displayed code has its stated
original PTS-rk type. The type
$\mathsf{Ch}$ has rank $2$, and the three operation types have rank $3$.
The exact counts, arithmetic laws, and forward reduction of
$\overline n$ to the unapplied annotated normal form
$\lambda f:E.\lambda x:X.f^n x$ are checked.
%proof{rocq-local:BCKW.v}
:::

This establishes the numeral fragment, not completeness for arbitrary
dependent proofs. The BCKW QIIRT of Section <#sec-qiirt> likewise checks as
an independent signature. To replace the complete Hilbert catalogue one
must additionally derive its dependent substitution schemes, with their
literal annotations and types. That translation remains open.

# Why a fixed finite typed basis is impossible {#sec-finite-basis}

Write $\rho(t)=s$ when $t:T$ and $T:U_s$: this is the rank of the
*type* of $t$. In a well-formed context it is unique, invariant under
beta conversion between typable terms, and satisfies

$$
\rho(fu)\leq\rho(f).
$$

Every universe atom has type rank zero, since $U_i:U_0:U_0$.
For finitely many fixed closed typed atoms $b_1,\ldots,b_m$, set
$R=\max(0,\rho(b_1),\ldots,\rho(b_m))$. Induction on application trees
bounds every typable expansion by $R$, even when every $U_i$ is available.
The bound also holds for typable beta-convertible outputs. But define

$$
\begin{aligned}
A_0&=U_0,& A_{s+1}&=\Pi x:A_s.U_0,\\
v_0&=U_0,& v_{s+1}&=\lambda x:A_s.U_0.
\end{aligned}
$$

Then $A_s:U_s$ and $v_s:A_s$, so $v_{R+1}$ cannot be represented.

:::Theorem[No finite fixed typed basis]
No finite list of closed originally typed PTS-rk terms is applicatively
complete up to original beta conversion, even with all universe atoms
available. The representing application tree must itself be originally
typable. %proof{rocq-local:FiniteBasis.v}
:::

This excludes a fixed typed BCKW basis or a single fixed typed combinator,
but not numerically indexed families with unbounded type ranks. Such
families supply the economical presentation below.

# Dependent logical schemes {#sec-core}

Only the natural numbers in the subscripts are external indices. Every
displayed type or family parameter is an ordinary term argument. Define

$$
\begin{aligned}
I_a&=\lambda A:U_a.\lambda x:A.x,\\
K_{a,b}&=\lambda A:U_a.\lambda B:U_b.\lambda x:A.\lambda y:B.x,\\
P_{a,b}&=\lambda A:U_a.\lambda B:(\Pi x:A.U_b).\Pi x:A.Bx.
\end{aligned}
$$

Their axiom types are

$$
\begin{aligned}
I_a&:\Pi A:U_a.\,A\to A,\\
K_{a,b}&:\Pi A:U_a.\Pi B:U_b.\,A\to B\to A,\\
P_{a,b}&:\Pi A:U_a.\Pi B:(\Pi x:A.U_b).\,U_{r(a,b)}.
\end{aligned}
$$

For dependent substitution, abbreviate the parameter telescope by

$$
\Theta_{a,b,c}=
\bigl(A:U_a,\ B:\Pi x:A.U_b,\ C:\Pi x:A.\Pi y:Bx.U_c\bigr).
$$

Here $\lambda\Theta$ and $\Pi\Theta$ bind its three entries in order.
The fourth scheme is

$$
\begin{aligned}
S_{a,b,c}
 &=\lambda\Theta_{a,b,c}.
   \lambda f:(\Pi x:A.\Pi y:Bx.Cxy).
   \lambda g:(\Pi x:A.Bx).\lambda x:A.fx(gx),\\
S_{a,b,c}
 &: \Pi\Theta_{a,b,c}.
    \Pi f:(\Pi x:A.\Pi y:Bx.Cxy).
    \Pi g:(\Pi x:A.Bx).\Pi x:A.Cx(gx).
\end{aligned}
$$

Logically, $I$ proves reflexive implication, $K$ adds an unused hypothesis,
and $S$ combines two deductions using the same hypothesis. The family $C$
allows the conclusion to depend on both that hypothesis and the intermediate
proof $gx$. The combinator $P$ forms the quantified proposition itself.

:::Theorem[Sound core schemes]
For every numerical index, the four displayed terms have exactly the
displayed original PTS-rk types. Their fully applied beta laws are

$$
\begin{aligned}
I_a\,A\,x&\longrightarrow_\beta^*x,&
K_{a,b}\,A\,B\,x\,y&\longrightarrow_\beta^*x,\\
S_{a,b,c}\,A\,B\,C\,f\,g\,x&\longrightarrow_\beta^*fx(gx),&
P_{a,b}\,A\,B&\longrightarrow_\beta^*\Pi x:A.Bx.
\end{aligned}
$$

The ranks of their types are respectively
$a+1$, $\max(a+1,b+1)$, $\max(a+2,b+2,c+1)$ and $a+2$.
%proof{rocq-local:Core.v}
:::

The rank of the type of $P_{a,b}$ is independent of $b$ because its result
is a universe symbol, itself of type $U_0$.

# A complete Hilbert presentation {#sec-hilbert}

Take the sorts $\{U_i:i\in\mathbb N\}$, the axioms $U_i:U_0$, and the
four preceding schemes. Complete the scheme set $\mathcal B$ by the following
operations of Bunder--Dekkers, Definition 7.3 [1]. They operate on schemes
with sort parameters, before numerical instantiation.

1. **Types of schemes.** From $M:T$ in $\mathcal B$, with $T$ not a sort,
   include $T':U_s$ whenever the original calculus types this scheme.
   Here $T'$ is $T$, or is obtained by replacing a repeated occurrence of
   a sort parameter in $T$ by a fresh sort parameter.

2. **Expose a product.** Whenever $\mathcal B$ contains

   $$
   (\lambda x_1:A_1\cdots\lambda x_{k-1}:A_{k-1}.
      \Pi x_k:A_k.B)
   : (\Pi x_1:A_1\cdots\Pi x_{k-1}:A_{k-1}.U_s),
   $$

   also include the corresponding well-typed scheme

   $$
   (\lambda x_1:A_1\cdots\lambda x_k:A_k.B)
   : (\Pi x_1:A_1\cdots\Pi x_k:A_k.U_t).
   $$

Only the sort constraints needed for each new scheme's own typing are
retained. They are constraints on fixed scheme shapes, not assumptions
that the theorem being translated is already derivable. The resulting
catalogue has 11 product, 7 identity, 17 constant and 30 substitution shapes
[1, Section 7]. Specialize their sort constraints using $U_i:U_0$ and
$r(a,b)=\max(a+1,b)$. Thus one family indexed by a tag from 65 shapes
and a finite tuple of natural numbers suffices. The number 65 is an upper
bound from this construction, not a minimality result.

The judgment $\vdash_H M:A$ has empty context. Its leaves are the universe
axioms and instances of $\mathcal B$; its inference rules are

$$
\begin{gathered}
\frac{\vdash_H M:\Pi x:A.B\quad\vdash_H N:A}
 {\vdash_H MN:B[N/x]}
\qquad
\frac{\vdash_H M:A\quad\vdash_H B:U_s\quad A=_\beta B}
 {\vdash_H M:B},\\[1em]
\frac{\vdash_H M:A\quad M\longrightarrow_\beta N}
 {\vdash_H N:A}
\qquad
\frac{\vdash_H M:A\quad A\longrightarrow_\beta B}
 {\vdash_H M:B}.
\end{gathered}
$$

These are application, conversion, subject reduction and type reduction.
There is no abstraction or product-formation inference rule: their effect
is recovered from the schemes. Annotated lambda and product syntax remains
inside the schemes, types and reductions.

PTS-rk is *supersorted*: every constant has a sort and every pair of sorts
has a product rule. On well-formed contexts, type correctness and
supersortedness identify the repository's sorted-domain abstraction rule
with the standard sorted-product rule [1, Theorem 7.2]. Hence
Bunder--Dekkers, Theorem 9.4, gives

$$
\boxed{\ \vdash_H M:A\quad\Longleftrightarrow\quad
             \vdash_{\rm rk}M:A\ }.
$$

This is an external completeness theorem, not a Rocq theorem or axiom of
this development. It supplies a Hilbert logical system: closed
proofs of exactly the original closed judgments arise from fixed schematic
logical axioms and the four rules above. Only its annotated, beta-only
presentation is used here.

:::Theorem[Hilbert soundness and proof extraction]
For any catalogue whose original atom typings have been independently
proved, a Hilbert derivation of $M:A$ yields a binder-free applicative code
$c$ such that

$$
\vdash_{\rm rk}M:A,\qquad
\vdash_{\rm rk}E(c):A,\qquad E(c)\longrightarrow_\beta^*M.
$$

Here $E$ expands catalogue atoms and preserves application.
%proof{rocq-local:Hilbert.v}
:::

The atom-typing premise belongs to the metatheorem, not the definition of a
Hilbert derivation. Extraction retains a code through both reduction rules;
its expansion recovers the subject by forward beta reduction, not merely
conversion. For the literal 65-shape catalogue below,
atom typing is discharged: soundness and extraction hold without an
additional premise. %proof{rocq-local:CatalogueEconomy.v}
The extracted code represents a proof through the original typing of
$E(c)$; types still use original annotated syntax.

## The literal catalogue and its completeness boundary {#sec-catalogue}

A literal catalogue can be described by the four core schemes, their types,
and finite prefixes of their parameter telescopes. The types of the core
schemes and their successive product exposures are included. For each
prefix of length $k$ listed below, put an arbitrary universe $U_d$ at its
end and replace the first $e$ products by lambdas, for every
$0\leq e\leq k$:

| Telescope | Prefix lengths | Universe-tail templates, in exposure order |
| --- | --- | --- |
| $P$ | $2;3$ | $P_2,P_3,P_8$; $P_7,P_9,P_{10},P_{11}$ |
| $I$ | $2$ | $I_5,I_6,I_7$ |
| $K$ | $2;3;4$ | $K_{12}$--$K_{14}$; $K_{15}$--$K_{18}$; $K_7$--$K_{11}$ |
| $S$ | $3;4;5;6$ | $S_9$--$S_{12}$; $S_{13}$--$S_{17}$; $S_{18}$--$S_{23}$; $S_{24}$--$S_{30}$ |

Length zero gives a universe; length one gives
$P_5=\Pi A:U_a.U_d$ and $P_4=\lambda A:U_a.U_d$.
The $S$ prefix of length two is the $P$ prefix of length two.
The length-three $P$ telescope is
$(A:U_a,B:\Pi x:A.U_b,x:A)$; its terminal body $Bx$ supplies
$P_6=\lambda A:U_a.\lambda B:(\Pi x:A.U_b).\lambda x:A.Bx$.
The other telescopes are the full binder lists of $I,K,S$ above.
All domain-sort slots and the final slot $d$ are independent.
The labels are local names, not an identification with the paper's numbering.

These templates give 11 product, 7 identity, 18 constant and 30 substitution
shapes. Every instance has a proved original typing and exact type rank.
%proof{rocq-local:ProductCatalogue.v}
%proof{rocq-local:IdentityCatalogue.v}
%proof{rocq-local:ConstantCatalogue.v}
%proof{rocq-local:SubstitutionCatalogue.v}
Every declared type is literally another catalogue subject or a universe;
its Hilbert type judgment and every leading-product exposure are derivable.
%proof{rocq-local:CatalogueHilbert.v}
The redundant constant shape
$K_{13}(a,b,c)=\lambda A:U_a.\Pi B:U_b.U_c$ has the independently typed
representation

$$
K_{0,0}\,U_1\,U_a\,(P_5(b,c))
 \longrightarrow_\beta^*\lambda A:U_a.\Pi B:U_b.U_c.
$$

Removing it leaves 65 shapes and exactly the same Hilbert judgments.
%proof{rocq-local:CatalogueEconomy.v}
The remaining $P_4,P_5$ also have independently typed representations in
the four-family core, with forward beta recovery.
%proof{rocq-local:ProductCatalogue.v}

:::Corollary[Completeness of the literal catalogue]
The 65-shape Hilbert system derives exactly the closed judgments of
original PTS-rk. This mathematical result uses [1]; its completeness
direction is not yet a Rocq theorem.
:::

To prove it, first retain all 66 templates. A declared type is a universe,
a core type, or a pure-product prefix ending in a universe. Freshening a
repeated sort occurrence in operation (1) therefore only renames an
independent domain slot or frees the final slot; the table contains every
well-typed instance. Operation (2) increments $e$, and all exposures are
present. The original functional sort rules fix each declared type.
Thus the 66 templates contain the schematic catalogue of Definition 7.3;
Theorem 9.4 supplies completeness by inclusion. Soundness supplies the
reverse implication, and the checked 65/66 equivalence removes $K_{13}$.
The symbolic template-inclusion step is not formalized in Rocq.

The formalized support includes contextual start and empty-context
equivalence, %proof{rocq-local:ContextualHilbert.v}
internal substitution, %proof{rocq-local:ContextualSubstitution.v}
classification of the seven variable/application-ended heads,
%proof{rocq-local:CatalogueShapes.v}
literal product-exposure compatibility and beta simulation,
%proof{rocq-local:ProductExposure.v}
$I,K,S$ abstraction and product assembly given the required family
derivations, %proof{rocq-local:HilbertAbstraction.v}
and the local type-preserving six-to-three $S$ contraction with shared
premises. %proof{rocq-local:HilbertReplay.v}
General chain shortening with a sharing-aware measure, internal type
correctness and unconditional deduction admissibility remain unformalized.
The badges certify the supporting results, not the completeness direction.

## A Metamath Zero specification {#sec-mmzero}

This Hilbert calculus has a first-order specification in Metamath Zero
(MM0) [6].
The accompanying [CombinatoryPTS-rk.mm0](CombinatoryPTS-rk.mm0) uses strict
sorts for natural-number indices, raw de Bruijn terms, finite ranked
telescopes and provable statements. Lambda annotations remain part of the
term syntax; beta reduction retains capture-avoiding substitution.
Auxiliary Horn rules compute ranks, shifts, substitutions and reductions.
The graph `form e L D E M T` folds a ranked telescope `L` over a body/type
pair `D:E`, exposing `e` initial lambdas and leaving products afterwards.
Only fixed axiom rows feed Hilbert typing: they generate the 11 product,
7 identity, 17 constant and 30 substitution shapes, without an opaque
catalogue lookup or original-typing oracle.

Its Hilbert inference layer is the following literal source excerpt:

%CombinatoryHilbertListing{}

Here an MM0 `axiom` declares a rule of this specified logical system;
the arrow between dollar-delimited statements separates premises from
conclusion. It is not a new assumption in the Rocq development. Read the
rules by their least syntactic derivability relation, not as a claim that
every first-order model has only the intended numbers or terms. Only the
fixed closed templates and universe terms initialize Hilbert typing;
application, conversion and reduction then preserve closedness.
In `h_convert`, `C` is a common beta reduct of the two types; original
confluence identifies this condition with beta conversion.
An MM0 proof certifies such a derivation, while the PTS subject is the proof
object it types. The specification and
[ten example derivations](../../docs/direct-r-handoff/verification/combinatory-mm0-2026-10-22/CombinatoryPTS-rk-examples.mm1),
including capture-avoiding substitution, pass the MM0 elaborator and
reference checker.
These checks do not establish the general Rocq correspondence or source
completeness.
The full source is included in Appendix <#app-mmzero>, directly from the
same file; it is longer than one page.

# A first-order algebraic presentation as QIIRT {#sec-qiirt}

A generalized algebraic theory (GAT) specifies dependent sorts, operations
and equations; its initial algebra is syntax modulo those equations [3, 4].
Here we present the contextual BCKW(I) signature with a product former,
separately from the full Hilbert catalogue.
The natural-model presentation [5] keeps the context as an index and computes
each term's type recursively. The result is a quotient
inductive-inductive-recursive type (QIIRT): its sorts, paths and typing
function are defined together. The following excerpts come directly from
the checked Cubical Agda source, [BCKW.agda](agda/BCKW.agda).

%begin{minipage}{\linewidth}
%CombinatoryAgdaListing{agda-signature}
%end{minipage}

Here $\operatorname{Ty}(\Gamma,r)$ represents types classified by $U_r$.
Thus $\operatorname{RankedTy}\Gamma=\sum_r\operatorname{Ty}(\Gamma,r)$
retains the sort of the type, and the fibre

$$
\operatorname{TmAt}(\Gamma,r,A)
 =\sum_{t:\operatorname{Tm}\Gamma}(\operatorname{tyOf}t=(r,A))
$$

recovers terms at a specified type. Equality constraints serve in place of
a type index: this is the Ford transformation, not untyped syntax.

## Operations and equations

The context operations are empty context and extension
$\Gamma\mathbin{\triangleright}A$. Substitutions have identity, composition,
the unique substitution to the empty context, guarded pairing, and
projections. Write their actions as $A[\sigma]$ and $t[\sigma]$.
In addition to the category laws, impose

$$
\begin{gathered}
A[\mathrm{id}]=A,\qquad
A[\sigma][\nu]=A[\sigma\circ\nu],\\
t[\mathrm{id}]=t,\qquad
t[\sigma][\nu]=t[\sigma\circ\nu],\\
\operatorname{fst}(\sigma,t)=\sigma,\quad
\operatorname{snd}(\sigma,t)=t,\quad
(\operatorname{fst}\sigma,\operatorname{snd}\sigma)=\sigma,\\
(\sigma,t)\circ\nu=(\sigma\circ\nu,t[\nu]).
\end{gathered}
$$

Pairing requires $\operatorname{tyOf}t=(r,A[\sigma])$.
The singleton $\langle t\rangle=(\mathrm{id},t)$ uses the type-identity
path; lifting a substitution through a binder uses the type-composition
path. These structural guards are proved in the source.
Comprehension eta is a law of substitutions, not function eta.

The logical operations are universe types $U_i:\operatorname{Ty}(\Gamma,0)$,
universe terms $\mathrm{univ}_i:\operatorname{Tm}\Gamma$, guarded decoding
$\mathrm{El}_r$, dependent products and application. Their types satisfy

$$
\begin{gathered}
\operatorname{tyOf}(\mathrm{univ}_i)=(0,U_0),\qquad
\mathrm{El}_0(\mathrm{univ}_i)=U_i,\\
\frac{A:\operatorname{Ty}(\Gamma,a)\quad
 B:\operatorname{Ty}(\Gamma\triangleright A,b)}
 {\Pi A B:\operatorname{Ty}(\Gamma,\max(a+1,b))}.
\end{gathered}
$$

A codomain is syntax in an extended context, not a meta-function on terms.
Application is exported with the usual argument-type guard:

%CombinatoryAgdaListing{agda-application}

The primitive guard uses $A[\mathrm{id}]$ to keep recursive definitions
strictly positive; the displayed wrapper derives it from the identity law.
Substitution is natural for universes, decoding and products, with
$(\Pi A B)[\sigma]=\Pi(A[\sigma])(B[\sigma\uparrow_A])$.
The closed atoms are $B_{a,b,c},C_{a,b,c},K_{a,b},W_{a,b},I_a$
and $P_{a,b}$. The first five explicitly quantify over their type
arguments at the ranks indicated by the subscripts, then have the types
of Section <#sec-bckw>; $P$ has the type of Section <#sec-core>.
Their literal telescopes are defined directly in the source.
Typed application spines fix the saturation arities of the
guarded $B,C,K,W,I$ laws. The $P$ law is
$\mathrm{El}_{r(a,b)}(P_{a,b}AB)=\Pi x:\mathrm{El}_a(A).\mathrm{El}_b(Bx)$,
with the displayed applications and decodings well typed.
There is no abstraction constructor, function eta or unrestricted quotation
$\operatorname{Ty}(\Gamma,r)\to\operatorname{TmAt}(\Gamma,0,U_r)$.

For a term path $L=R$, recursion must also supply a path between its types.
The guarded path constructors therefore carry
$e:\operatorname{tyOf}L=\operatorname{tyOf}R$, with recursive clause
$\operatorname{tyOf}(\mathrm{law}\ e\ i)=e(i)$.
Only the specified laws generate paths, not arbitrary pairs of equally
typed terms. All four sorts are sets: the term-truncation coherence guard
is derived from the sethood of $\operatorname{RankedTy}\Gamma$.
Structural contraction already follows from comprehension:
in $\Gamma\triangleright A$, the substitution
$(\mathrm{id},q_A)$ maps into
$\Gamma\triangleright A\triangleright A[p]$. Substituting along it shares
one variable. This is distinct from compiling every dependent term into
closed $B,C,K,W,I$ atoms; ordinary contextual application itself permits
shared variables.

This gives executable, intrinsically typed quotient syntax. The source
derives the structural guards; general admissibility of the logical
equations' coherence guards, interpretation in original annotated PTS-rk,
equality reflection and dependent combinatory completeness remain separate
obligations.

# Elementary normalization and finite syntax carriers {#sec-finite-normalization}

## The rank parameter and the reduction strategy

Let $|t|$ count all constructors of the original annotated syntax, including
annotations. Use Knuth's tetration notation
$a\uparrow\uparrow0=1$ and
$a\uparrow\uparrow(h+1)=a^{a\uparrow\uparrow h}$, and put $q=\max(2,N)$.
The relevant bound $d$ bounds the type rank of **every redex head** in $t$,
also inside annotations. It is the repository's `RedexBound`, not merely
the outer type rank $\rho(t)$. Indeed the closed terms

$$
t_a=(\lambda x:(\Pi A:U_a.A\to A).U_0)I_a
$$

all have type $U_0$ and $13$ nodes, while their root redex has rank $a+2$.
This separates the two rank notions; it is not a lower bound on the cost
of reducing these particular terms.

:::Theorem[Normal form and intermediate-size bound]
Suppose $\Gamma$ is well formed, $\Gamma\vdash t:A$, $|t|=N$, and
`RedexBound` holds at $d$. There is a deterministic beta-reduction sequence
to a normal $t^*$ with $\Gamma\vdash t^*:A$. Every intermediate term,
including $t^*$, has at most $q\uparrow\uparrow(d+1)$ nodes and retains
redex-rank bound $d$.
The documented abstract normalization cost is at most
$(N+5d+2)\uparrow\uparrow(2d+2)$.
%proof{rocq-local:QuantitativeNormalization.v}
:::

Use $d$ passes of the existing greedy normalizer [8]. Within a pass, reduce
the annotation or left child first, then the body or right child, then
contract the root when possible. A pass lowers the maximal redex rank;
sequentialization bounds every single-beta intermediate by $2^{|t|}$.
Iterating this binary exponential $d$ times gives a bound no larger than
$q\uparrow\uparrow(d+1)$; the comparison is checked.
%proof{rocq-local:KnuthBounds.v}
Original subject reduction and redex-rank preservation supply the typing
and rank invariants. For fixed $d$, these functions are elementary, with
tower heights linear in $d$.

For any fixed finite set of atom shapes, including the indexed catalogue,
expansion satisfies $|E(c)|\leq C\ell$ when $c$ has $\ell$ atoms, with a
shape-dependent constant $C$. Thus
$\max(2,C\ell)\uparrow\uparrow(d+1)$ bounds normalization of its expansion.
For type-indexed atoms the bound requires fixed type shapes, as in the
numeral family; arbitrary type-index payloads must also be charged.
This does not yet bound the length of a normal **binder-free BCKW output**:
an annotation-preserving compiler back to that syntax, with a size theorem,
is still needed. Numerical labels count as one node here;
their bit lengths must be charged in a machine-level complexity claim.

## Finite normal families and dependent application

Rank and length alone do not make normal forms finite: the infinitely many
distinct $U_i:U_0$ are all closed normal terms of size one and type rank zero.
Fix instead a finite alphabet $\Sigma$ of universe labels. For a fixed type
$A$, the desired family of **syntax codes** is

$$
\mathrm{NF}_{\Sigma}(\ell,r;A)=
\{t\mid\operatorname{labels}(t)\subseteq\Sigma,
 |t|=\ell,\ t\text{ normal},\ \vdash t:A,\ \vdash A:U_r\}.
$$

This is a proposition-valued fibre of finite syntax, not a collection
counting arbitrary typing derivations. The checked Agda construction below
uses genuine normal representatives and an injective finite encoding.
Universe labels and de Bruijn variables range over finite alphabets;
products and lambdas retain both their annotation and body.

:::Theorem[Finite carriers]
Beta reduction introduces no universe labels. All closed terms of at most
$s$ nodes over $\Sigma$, and hence all their normal subfamilies, have an
explicit complete finite enumeration of length at most
$(3s+|\Sigma|+4)\uparrow\uparrow3$ (duplicates are harmless). Consequently, for a closed typed
$d$-bounded input of size $N$, every state of the preceding strategy lies
in a finite carrier of size at most

$$
(4N+4)\uparrow\uparrow(d+3),\qquad \Sigma=\operatorname{labels}(t).
$$

For closed normal $f:\Pi x:A.B$ and $y:A$, with
$\Pi x:A.B:U_r$, let $N=1+|f|+|y|$. Running $r$ passes on $fy$ returns a
normal $z:B[y/x]$. Every intermediate has size at most
$\max(2,N)\uparrow\uparrow(r+1)$ and lies
in the finite raw carrier over $\operatorname{labels}(f)\cup
\operatorname{labels}(y)$ of capacity at most $(4N+4)\uparrow\uparrow(r+3)$.
%proof{rocq-local:FiniteNormalForms.v}
:::

The last rank bound needs no internal rank inference: normal $f,y$ have
no internal redexes, so only the new root redex matters. Its rank is $r$.
The algorithm can represent each term state by an index into the displayed
finite list; the proofs check the underlying normalizer and its finite
envelope, not an extracted finite-memory machine. They do not bound the
storage of arbitrary typing witnesses. Nor is an externally supplied,
possibly beta-expanded $B$ bounded by $|f|+|y|$: count supplied type syntax
and its rank budget separately if types themselves are to be normalized.

## Inductive-inductive-recursive normal representatives

The quotient syntax of Section <#sec-qiirt> identifies beta-equivalent
terms of different lengths. Length must therefore be measured on explicit
representatives. [NormalForms.agda](agda/NormalForms.agda) defines normal
syntax, a non-lambda-head family, and their raw interpretation together:

%begin{minipage}{\linewidth}
%CombinatoryAgdaListing{agda-normal-iir}
%end{minipage}

The occurrence of `erase` in the `non-lam` constructor makes this a genuine
inductive-inductive-recursive definition. Here `n` is context length, while
`erase` retains every original annotation; it forgets only the normal-form
presentation. `NotLambda` excludes precisely lambda heads. Thus application
cannot introduce a beta redex, and all children, including annotations,
are normal. Conversely, every normal raw term has such a representative.
The head witnesses are propositions and `erase` is injective.
%proof{agda-local:NormalForms.agda}

Exact and bounded node counts give the fibres

%begin{minipage}{\linewidth}
%CombinatoryAgdaListing{agda-normal-family}
%end{minipage}

An `AtMost A k` certificate is a function `A → Fin k` together with its
injectivity proof. Depth-bounded raw syntax has a sum/product encoding
equivalent to `Fin (capacity a n s)`; every node-bounded term encodes and
decodes identically. The arithmetic proof bounds this capacity by the
displayed tetration. No normalization theorem is assumed.
%proof{agda-local:FiniteCarriers.agda}
%proof{agda-local:FiniteBounds.agda}

[NativeTyping.agda](agda/NativeTyping.agda) transcribes the original six
typing rules and annotated beta equivalence, not a bidirectional checker.
For a fixed closed type `A`, `NormalAt labels N d A t` asserts normality,
exactly `N` nodes, original typing at `A`, and a sort rank of `A` at most
`d`. Typing and rank certificates are propositionally truncated so that
different derivations do not create additional inhabitants. The resulting
typed family and its cardinal bound are checked directly:

%CombinatoryAgdaListing{agda-typed-normal}

The constructor `pack-typed-normal` includes every originally typed normal
term over the chosen alphabet. Here $a\le N$ is essential for the
alphabet-independent numerical bound: take the labels occurring in an
input of size `N`, not an arbitrarily large ambient alphabet. The dictionary
`labels : Fin a → Nat` and type `A` are fixed parameters. Duplicate dictionary
entries merely overcount representatives; they do not invalidate the bound.
The parameter `d` bounds the type rank in this normal family; a normal
term's redex bound alone would be vacuous.
%proof{agda-local:TypedNormalFibres.agda}

## Total typed apply-and-normalize

[BoundedReduction.agda](agda/BoundedReduction.agda) implements capture-avoiding
substitution, the greedy pass and its iteration on original annotated syntax.
It proves forward beta reduction, decidable normality, and compatibility
with replacing finite label indices by their original universe numbers.
The size proof is unconditional: one pass has at most $2^{|t|}$ nodes,
and every single-beta intermediate of $d$ passes has at most $T_d(|t|)$,
where $T_0(N)=N$ and $T_{d+1}(N)=2^{T_d(N)}$.
This seeded binary tower is bounded by
$\max(2,N)\uparrow\uparrow(d+1)$.
%proof{agda-local:RawSizeBounds.agda}
%proof{agda-local:BoundedTrace.agda}

In [FiniteTrace.agda](agda/FiniteTrace.agda), a state is literally
`State a n s = Fin (capacity a n s)`. Decoding recovers the annotated term;
`state-bound` injects the closed state type, for $s=T_d(N)$, into
`Fin` of size $(4N+4)\uparrow\uparrow(d+3)$. The computed beta path has
finite-indexed states and proposition-valued edge evidence:

%CombinatoryAgdaListing{agda-finite-walk}

Here `IndexedStep i j` is the propositional truncation of an original
beta step between the decoded indices. `rounds-trace` computes this walk
with the exact input and output endpoints; `walk-sound` proves forward
reduction after decoding. Every state has the bound, not the type of
arbitrarily long whole walks. %proof{agda-local:FiniteTrace.agda}

For two closed IIR normal forms, put `N = 1 + |f| + |y|`.
[FiniteApplication.agda](agda/FiniteApplication.agda) defines the successful
result type and proves its entire cardinal bound, including proof-irrelevant
reduction evidence:

%CombinatoryAgdaListing{agda-finite-result}

The raw function `applyWithinFinite` returns `Maybe ApplicationResult`:
without typing premises, the selected number of passes need not normalize.
Its successful result is an actual normal representative with a certified
bounded forward reduction. Size cannot cause failure.
%proof{agda-local:FiniteApplication.agda}

[TypedApplication.agda](agda/TypedApplication.agda) packages the original
premises $f:\Pi x:X.Y$, $y:X$, and $\Pi x:X.Y:U_d$ in `TypedInputs`.
The inputs $f,y$ are IIR normal forms; no normalization or output-typing
certificate is an input. `TypedResult` includes the exact original output
typing at $Y[y/x]$, the bounded forward reduction, and the same cardinal
bound. The executable interface is total:

%begin{minipage}{\linewidth}
%CombinatoryAgdaListing{agda-typed-apply}
%end{minipage}

`normal-result` constructs the normal representative of `rounds d (app f y)`;
`output-typed` certifies its original dependent type. Typing premises are
propositionally truncated: they justify normality and output typing, but do
not select the computed syntax. The entire result type injects into
`Fin` of size $(4N+4)\uparrow\uparrow(d+3)$, with $N=1+|f|+|y|$.
%proof{agda-local:TypedApplication.agda}

The Agda proof uses the original six rules throughout. Binder-aware algebra
gives weakening and typed substitution, including annotations.
%proof{agda-local:NativeSubstitution.agda}
Parallel reduction proves confluence and uniqueness of normal forms.
%proof{agda-local:NativeConfluence.agda}
Typing inversion, context conversion and substitution prove type correctness
and subject reduction. %proof{agda-local:NativePreservation.agda}
Type-rank uniqueness and substitution safety show that one greedy pass
lowers a positive redex bound; bound zero implies normality. Hence $d$
passes normalize, and normal arguments give the required application bound.
%proof{agda-local:NativeDescent.agda}
Every single-beta intermediate also preserves the original rank bound.
%proof{agda-local:NativeRankPreservation.agda}
`normalize` packages these results for any originally typed, rank-bounded
term, including its normal form, exact typing and bounded trace. It is
idempotent and independent of the sufficient rank budget; two normalized
inputs are beta-convertible exactly when their normal outputs are equal.
%proof{agda-local:NativeNormalization.agda}

[NativeRankCoverage.agda](agda/NativeRankCoverage.agda) computes a sufficient
redex bound from every original typing derivation in a well-formed context,
recursing through the annotated subject and using type-rank uniqueness.
Thus `normalize-typed` requires neither a supplied rank bound nor a
termination certificate.

The finite-indexed trace certifies the syntax-state bound. Its construction
still evaluates unrestricted `Raw` internally; neither it nor the result
bound asserts that every host type, proof object or runtime allocation is
bounded. The following finite-table compilation makes the executable state
interface explicit without claiming such a whole-runtime bound.

## Church compilation at explicit ranks {#sec-church-finite}

Fix a nonempty label alphabet of size $a\le N$, and put
$s=T_d(N)$ and $K=\operatorname{capacity}(a,0,s)$, so
$1\le K\le(4N+4)\uparrow\uparrow(d+3)$.
For every positive $m$ and fixed rank $r$, define the original PTS terms

$$
\begin{aligned}
\mathsf E_m^r
 &=\Pi X:U_r.\;\Pi x_0:X.\cdots\Pi x_{m-1}:X.\;X,\\
\mathsf c_i^{m,r}
 &=\lambda X:U_r.\;\lambda x_0:X.\cdots\lambda x_{m-1}:X.\;x_i
 \qquad(0\le i<m).
\end{aligned}
$$

Then $\mathsf E_m^r:U_{r+1}$ and
$\mathsf c_i^{m,r}:\mathsf E_m^r$. Every inner arrow has rank $r+1$;
the outer product has rank $\max(1,r+1)=r+1$ because $U_r:U_0$.
Both the type and its canonical selectors are normal and have exactly
$2m+3$ annotated nodes. These facts, including original typing for every
index, are checked in [ChurchFinite.agda](agda/ChurchFinite.agda).
%proof{agda-local:ChurchFinite.agda}

The finite injections already proved give the explicit encoding

$$
\ulcorner t\urcorner_r=\mathsf c_{\iota(t)}^{K,r},
\qquad
\iota(t)=\operatorname{encodeState}(t).
$$

For the exact-length normal family one can use the smaller
$K_N=\operatorname{capacity}(a,0,N)$ instead. Typing/rank proofs select an
external subfamily of these codes; they are not additional stored data.
This encodes the represented elements of the Agda fibres, not their Cubical
paths or an internal dependent induction principle.

For a finite table $F:\mathrm{Fin}(m)\to\mathrm{Fin}(n)$, define

$$
\begin{aligned}
\mathsf{Map}_r(F)
 &=\lambda p:\mathsf E_m^r.\;
   \lambda X:U_r.\lambda y_0:X.\cdots\lambda y_{n-1}:X.\;
   p\,X\,y_{F(0)}\cdots y_{F(m-1)},\\
\mathsf{Map}_r(F)&:\mathsf E_m^r\to\mathsf E_n^r,\qquad
(\mathsf E_m^r\to\mathsf E_n^r):U_{r+2},\\
\mathsf{Map}_r(F)\,\mathsf c_i^{m,r}
 &\longrightarrow_\beta^*\mathsf c_{F(i)}^{n,r}.
\end{aligned}
$$

This uses the same answer type $X$ throughout: it does not instantiate
$p$ at the higher-ranked type $\mathsf E_n^r$, and requires no cumulativity.
The original typing, displayed ranks and forward beta computation of
selectors, unary maps and their finite iterations are checked directly
against the native PTS rules.
%proof{rocq-local:ChurchFinite.v}

For a binary table $F(i,j)$, abbreviate its rows by
$R_i=q\,X\,y_{F(i,0)}\cdots y_{F(i,n-1)}$ and use

$$
\begin{aligned}
\mathsf{Map}^{(2)}_r(F)
 &=\lambda p:\mathsf E_m^r.\lambda q:\mathsf E_n^r.\lambda X:U_r.\\
 &\qquad\lambda y_0:X.\cdots\lambda y_{\ell-1}:X.\;
 p\,X\,R_0\cdots R_{m-1}.
\end{aligned}
$$

Its type is $\mathsf E_m^r\to\mathsf E_n^r\to\mathsf E_\ell^r$,
again of rank $r+2$, and it computes $F(i,j)$ on canonical inputs.
Indeed each row has type $X$, and successive elimination of the two
selectors chooses $y_{F(i,j)}$. This binary generalization is implemented
and regression-tested in Agda; its all-index typing and computation
theorems are not yet formalized.
Consequently rank $1$ suffices for the finite data and rank $2$ for the
operators, taking $r=0$. The number of branches, hence program and type
size, grows with the external bounds; this is not a uniform fixed-interface
encoding of all elementary computations [8].

### Compiling reduction passes

Define a total finite transition $g:\mathrm{Fin}(K)\to\mathrm{Fin}(K)$
by decoding a state, performing one greedy pass, and re-encoding when the
result has at most $s$ nodes; otherwise leave the state unchanged.
Similarly let $b(i,j)$ encode the application of the decoded terms when it
fits, and return $i$ otherwise. These fallback branches make the tables
total on the entire ambient carrier. They are never used on the intended
bounded runs.

[FiniteMachine.agda](agda/FiniteMachine.agda) implements these functions.
Writing $t_{ij}=\operatorname{app}(\operatorname{decode}i,\operatorname{decode}j)$,
it proves, without a typing premise,

$$
\begin{gathered}
|t|\le N
\quad\Longrightarrow\quad
\operatorname{decode}(g^d(\iota(t)))=\operatorname{rounds}(d,t),\\
T_d(|t_{ij}|)\le s
\quad\Longrightarrow\quad
\operatorname{decode}(g^d(b(i,j)))=
\operatorname{rounds}(d,t_{ij}).
\end{gathered}
$$

%proof{agda-local:FiniteMachine.agda}

[FragmentMachineSoundness.agda](agda/FragmentMachineSoundness.agda) proves
typed soundness of these same finite-state functions. `finite-normalize`
decodes a normal form with the original type and rank bound; `finite-apply`
does so at exactly $Y[y/x]$ for normal typed inputs. Both return forward
reduction and bounded single-beta trace certificates. Their numerical
budget premises exclude clipping, not assume normalization.
%proof{agda-local:FragmentMachineSoundness.agda}

The original PTS programs are now explicit finite terms:

$$
\begin{aligned}
G&=\mathsf{Map}_r(g),&
A&=\mathsf{Map}^{(2)}_r(b),\\
\mathsf{Norm}_{a,d,N}^r
 &=\lambda p:\mathsf E_K^r.\;G^d p
 &: \mathsf E_K^r\to\mathsf E_K^r,\\
\mathsf{Apply}_{a,d,N}^r
 &=\lambda p:\mathsf E_K^r.\lambda q:\mathsf E_K^r.\;G^d(Apq)
 &: \mathsf E_K^r\to\mathsf E_K^r\to\mathsf E_K^r .
\end{aligned}
$$

Here $G^d p$ is a metalevel abbreviation for exactly $d$ nested applications,
not an added recursion primitive. The compiler tabulates one pass and the
application constructor, then emits these compositions; it does not fill a
table with precomputed final normal forms. For a closed originally typed
input $t$ over the chosen alphabet, with $|t|\le N$ and
$\operatorname{RedexBound}(\varepsilon,t,d)$, the decoded output is its
normal form. For normal
$f:\Pi x:X.Y$ and $y:X$, take $d$ as the product's sort rank and
$N=1+|f|+|y|$: the output decodes to a normal $z:Y[y/x]$ by the existing
normalization and preservation theorems.

The following executable generators return literal original syntax
(`Term 0`), not host functions standing in for PTS terms.
Here `compileMap` and `compileMap2` implement the displayed
tables, and `repeatMap` emits the $d$ applications:

%CombinatoryAgdaListing{agda-church-machine}

This is code-to-code computation, distinct from the total typed Agda
function above. It does not itself return an inhabitant of $Y[y/x]$ inside
PTS or translate arbitrary Agda programs and quotient/truncation proofs.
The generic native map proof and the Agda generators are separate
formalizations, not a checked cross-language compiler-correctness bridge.
Tabulation can be extremely large; no practical construction time or
whole-machine memory bound is asserted. In particular, the source budget
$N$ does not bound the generated program or its execution: even one
canonical state code has $2K+3$ nodes.

## Graded fragment soundness {#sec-fragment-soundness}

For a fixed label dictionary, let $\mathcal F_s(A)$ contain closed IIR normal
forms of at most $s$ annotated nodes with propositional original typing at
$A$. [FragmentEmbeddings.agda](agda/FragmentEmbeddings.agda) gives an
injection into `Fin (capacity a 0 s)`. Increasing the size budget gives
injective embeddings satisfying identity and composition, preserving the
decoded term and its typing. For these normal representatives,

$$
\operatorname{decode}(u)=_\beta\operatorname{decode}(v)
\quad\Longleftrightarrow\quad
\operatorname{decode}(u)=\operatorname{decode}(v).
$$

The right-hand side is literal equality of native annotated terms, not of
finite codes when the dictionary repeats labels. Original type conversion
acts on the typing fibre and commutes with the size embeddings.
%proof{agda-local:FragmentEmbeddings.agda}

Typed apply-and-normalize gives a graded application operation: for
$f\in\mathcal F_s(\Pi x:X.Y)$ and $y\in\mathcal F_t(X)$ its output belongs
to $\mathcal F_{T_d(1+s+t)}(Y[y/x])$, where $\Pi x:X.Y:U_d$.
The API retains the alphabet premise $a\le1+|f|+|y|$.
The output decodes to the unique normal form of $fy$;
subject reduction proves its exact type. Increasing budgets does not change
the decoded result, and beta-convertible inputs give beta-convertible
outputs. %proof{agda-local:FragmentApplication.agda}
These are sound closed-term operations and coherent finite inclusions, not
closure of one fixed carrier under every operation.

Alphabet inclusions preserve decoded terms and commute with normalization,
size embeddings, conversion and application. Injective alphabet maps give
injective carrier maps satisfying identity and composition.
%proof{agda-local:FragmentAlphabet.agda}
[AlphabetCoverage.agda](agda/AlphabetCoverage.agda) constructs finite prefix
alphabets for arbitrary annotated terms, contexts and valuations, with exact
decoding. A prefix alphabet's size can exceed the number of syntax nodes;
the general interpretation therefore includes alphabet size in its seed.

### Contextual certificate soundness

An independent finite proof-tree grammar records all six typing rules and
explicit beta instructions, including reduction inside annotations.
`Accepted` recursively checks subtrees and local syntax equalities;
`accepted?` decides this predicate without a native typing or normalization
oracle. Acceptance implies original `Typ`, and every original derivation
has an accepted certificate. This is checker adequacy, not a replacement
typing judgment. %proof{agda-local:NativeCertificates.agda}

Context certificates check each declaration's sort; valuation certificates
check every image against the substituted original context lookup. Their
decidable acceptance predicates derive well-formedness and typed
substitution. In particular, no native typing, well-formedness or
normalization premise is hidden in the following input predicate:
%proof{agda-local:CertificateContexts.agda}

%CombinatoryAgdaListing{agda-model-input}

`model-soundness` interprets every accepted input, not merely closed terms
or normal valuation images. It substitutes the environment into the subject
and type, computes a common sufficient rank bound, and normalizes both:

%CombinatoryAgdaListing{agda-model-soundness}

`ModelOutput` contains their finite normal representatives and forward
reductions, with original typing at the normalized type. Subject reduction
also retains exactly the substituted source type. The algorithm chooses
the alphabet and grades from the checked data; callers supply no guessed
budget or desired output. `check-and-interpret` returns `nothing` only for
rejected certificates, never because an accepted input failed to normalize.
%proof{agda-local:FragmentModelSoundness.agda}

The decoded interpretation respects conversion, is independent of typing witnesses
and sufficient budgets, and commutes with valuation composition:

$$
\operatorname{nf}\bigl(\operatorname{nf}(t[\sigma])[\tau]\bigr)
=\operatorname{nf}\bigl(t[\tau\circ\sigma]\bigr).
$$

Here valuations are arbitrary well-typed syntactic substitutions, not
arbitrary set-valued assignments. [NativeValuations.agda](agda/NativeValuations.agda)
proves identity, composition, context comprehension and substitution laws
for the original annotated syntax; normalization preserves those equations.
%proof{agda-local:ContextualNormalization.agda}

Finiteness counts the context and type as data, not as unbounded external
indices. [FragmentUniverses.agda](agda/FragmentUniverses.agda) bounds types,
values and valuations; [FragmentInterpretation.agda](agda/FragmentInterpretation.agda)
constructs their jointly normalized representatives. Write
$C_{a,n}(s)=\operatorname{capacity}(a,n,s)$ and
$D_{a,n}(s)=\prod_{j<n}C_{a,j}(s)$. The joint context/type/value package
$\mathcal P$ and context/substitution package $\mathcal M$ admit injections
with bounds

$$
\begin{aligned}
|\mathcal P_{a,n,s,d}|&\le D_{a,n}(s)\,C_{a,n}(s)^2,\\
|\mathcal M_{a,m,n,s}|&\le D_{a,m}(s)\,D_{a,n}(s)\,C_{a,n}(s)^m.
\end{aligned}
$$

Declarations retain their preceding scopes; all annotations and substitution
images are counted, while semantic certificates are propositional fibres.
Every typed valuation between well-formed contexts has such a package at
computed bounds. Images may remain raw: independently normalizing the
components of a dependent tuple is not assumed to preserve its exact type.
These bounds concern the represented packages, not arbitrary proof trees
or runtime allocations. %proof{agda-local:FragmentPackages.agda}

The interpretation is externally nontrivial. For
$\mathsf{Void}=\Pi X:U_0.X$, canonical-form inversion rules out a normal
inhabitant: a normal function must be a lambda, whose body would inhabit
the arbitrary type variable $X$. In the context $X:U_0$, no normal neutral
function or other normal constructor can do so. Total normalization gives
`Void-empty` for every original closed typing derivation, hence
`accepted-Void-rejected` and emptiness of every finite Void fibre.
These are Agda metatheorems, not internal PTS consistency proofs.
%proof{agda-local:NativeCanonicity.agda}

## Fragment interpretations and internal consistency {#sec-internal-consistency}

The result above is all-six-rule contextual soundness in a graded syntactic
interpretation. Every fixed typed input has computed finite coverage, and
larger grades cover more inputs. It is not yet a theorem that each stage is
a **finite first-order model of a logical fragment**. That requires a total
interpretation of the chosen signature on fixed finite domains, coherent
dependent fibres, and all chosen axioms under every semantic valuation.
Operations that increase grades do not automatically supply this closure.
For example, iterating $A_0=U_0$ and $A_{k+1}=\Pi x:U_0.A_k$ gives normal
types of rank one for $k>0$, using only label zero and $2k+1$ nodes: a fixed
size-bounded faithful type carrier cannot contain the whole family. This
does not rule out nonfaithful finite models.

The independent checker and its external soundness provide a precise
starting point for internal reflection. Their translation into PTS remains
open: equations on canonical Church codes do not give a finite-case or
dependent induction principle for arbitrary inhabitants of $\mathsf E_K^r$.
Compiler correctness must connect the encoded checker and interpretation
to the specified certificate predicate, including its conversion witnesses.

The intended endpoint is an original PTS proof of a specified consistency
sentence, not a metalevel assertion that no term of type $\mathsf{Void}$
exists. In the repository's [consistency plan](../self_consistency/README.tex),
this is $\vdash c:J(\mathcal{NAT}(C))$, where $C$ says that the fixed
predicate-only checker accepts no contradiction certificate, including all
its trace witnesses. A promising route uses supplied tower resources to
prove the checker's bounded soundness in elementary arithmetic, then
Pakhomov's transfer for predicate-only $\Pi_1$ sentences [9]:

$$
\mathsf{EA}\vdash C^{2\mathcal S}
\quad\Longrightarrow\quad
\mathsf H\vdash\mathcal{NAT}(C)
\quad\Longrightarrow\quad
\vdash c:J(\mathcal{NAT}(C)).
$$

Here $\mathcal S$ is the superexponential cut; the last implication requires
an actual interpretation $J$ of the stated theory into PTS. The finite
Church construction supplies computational ingredients, not those soundness
and interpretation proofs. External results for every fixed budget are not
one internal proof quantifying over all budgets. This distinction leaves
the Pakhomov-style self-consistency objective plausible without identifying
it with an unrestricted arithmetical consistency sentence. It also avoids
inferring consistency from the unrelated all-true Horn models below.

## A locally finitely satisfiable target {#sec-ground-horn}

There is a precise ground-query reduction, distinct from a faithful finite
model of syntax. Ground the Hilbert inference rules, including the auxiliary
syntax, substitution and reduction graphs. Introduce one nullary predicate
$p_J$ for each closed judgment $J$, and, for every ground rule instance, add

$$
J_1,\ldots,J_k\Rightarrow J
\quad\mapsto\quad
(p_{J_1}\land\cdots\land p_{J_k})\Rightarrow p_J.
$$

This is an effective positive Horn theory in a countably infinite recursive
vocabulary. Finite conversion certificates can be enumerated; no decision
procedure for untyped conversion is assumed. Grounding all the explicit
MM0 graph rules is a concrete version of the same construction.

:::Theorem[Ground derivability reduction]
A judgment $J$ has a finite source derivation exactly when the grounded
Horn theory entails $p_J$. Every finite fragment of the target has a
one-element first-order model.
%proof{rocq-local:GroundHorn.v}
:::

For preservation, induct on a derivation. For reflection, interpret $p_J$
as ``$J$ has a finite derivation'': this is the least model of the rules.
For finite satisfiability, use a singleton domain and make every $p_J$ true;
indeed this satisfies the whole positive theory. The badge proves the
generic Horn theorem; the effective Hilbert instantiation is the preceding
mathematical construction, not a checked MM0/native adequacy bridge.
Together with the external completeness theorem of Section <#sec-hilbert>,
it yields the stated ground-query reduction for original closed judgments.

This conclusion holds independently of the normalization bounds. Those
bounds additionally give finite representations of individual computations;
they do not supply a finite-signature interpretation preserving arbitrary
quantified formulas. Requiring finite models to contain faithful free
natural-number syntax would already fail: injective successor with zero
outside its image has no finite model. Thus neither a finite faithful
syntax model nor that stronger first-order reduction is claimed.

# The exact typed expansion quotient {#sec-exact-quotient}

Unlike the independent algebraic signature, this quotient is *relative to*
native annotated terms $\mathrm{Raw}$, applicative codes $\mathrm{Code}$,
expansion $E$, original typing $\mathrm{Typ}$ and beta equality $\mathrm{Beta}$;
$\mathrm{Uraw}(r)$ denotes $U_r$. The module [Exact.agda](agda/Exact.agda)
takes these data, not soundness or completeness proofs:

%begin{minipage}{\linewidth}
%CombinatoryAgdaListing{agda-exact}
%end{minipage}

Both endpoints of a term path are originally typed; type equality is
retained alongside equality of expansions. The common rank on each path
matches uniqueness of native sorts. The fibre of
$\operatorname{tyOfExact}$ recovers a specified ranked type without a term
type index. Given native beta reflexivity, the beta constructors identify
duplicate typing witnesses; set truncation makes equality witnesses irrelevant.

Ordinary inductive representatives with explicit typed beta relations are
also defined in the same file. Their type-respect lemma is checked; the
setoid laws follow when the supplied native beta relation is an equivalence.
Deleting quotient paths without retaining those relations would instead
give free syntax. Quotienting untyped codes first would be unsound as a
typing representation: a beta class can contain untypable representatives.

All accompanying Agda modules are checked with Agda 2.8.0, Cubical v0.9 and
the safe checker, without postulates, holes or disabled checks. This
establishes well-formedness of the constructions, not an Agda-to-Rocq
adequacy theorem or completeness of the BCKW signature.

# Future work {#sec-future-work}

The following three directions concern proposed systems, not changes to
the original calculus or results established above.

1. **A quadratically checkable PTS-rk variant.** Apply the ideas of van den
   Berg and den Besten [10]: replace implicit definitional conversion by
   explicit propositional computation and transport witnesses. Develop an
   objective variant preserving the rank discipline, and prove deterministic
   worst-case $O(L^2)$ type checking, where $L$ includes the context, term,
   type, numerical ranks and all annotations and witnesses. Their Theorem 7
   proves quadratic checking for objective type theory; adapting it to
   PTS-rk is a separate obligation. This bounds checking, not proof search
   or normalization.

2. **A logic-free arithmetic of partial functions.** Develop a system
   analogous to logic-free primitive recursive arithmetic, with the following
   requirements.

   - Terms denote partial functions, rather than necessarily total functions.
     The basic judgment is extension, not equality:

     $$
     f\sqsubseteq g
     \quad\Longleftrightarrow\quad
     \forall\vec x,y\,
       \bigl(f(\vec x)\mathbin{\downarrow}y
             \Rightarrow g(\vec x)\mathbin{\downarrow}y\bigr).
     $$

     Thus the graph of $f$ is included in the graph of $g$; symmetry is not
     an inference rule.
   - Every partial computable function must be expressible using a finite
     term-forming basis, in the spirit of finite bases for elementary
     recursive arithmetic. Computational completeness is a requirement on
     the intended standard interpretation, not on each finite model.
   - Use reflexivity, transitivity and substitution, together with a suitable
     induction rule. Substitution and induction must be sound for partiality
     and extension. Require a finite core of axioms and one infinite scheme
     $\mathsf{exists}_n$, whose intended meaning is: either all numbers are
     equal, or there is an initial segment of at least $n$ distinct natural
     numbers. This is a semantic specification to encode in the extension
     calculus, not an addition of unrestricted first-order logic.
   - Require local finite satisfiability: for every finite numeral $k$, the
     finite core together with all $\mathsf{exists}_n$ for $n\le k$ has a
     finite model validating the inference rules. Noncollapsed finite models
     realizing the retained initial segments are a stronger, useful goal;
     the collapse disjunct must not be mistaken for such a construction.
   - Require deterministic worst-case $O(L^2)$ derivation checking, or better,
     in the full encoded derivation length, including numerals, substitutions
     and induction certificates. Checking must not decide semantic extension
     of arbitrary programs or execute potentially divergent computations.

3. **Interpretation of the type theory in the arithmetic.** Construct an
   effective, derivation-preserving interpretation of the variant in (1)
   inside the system in (2). Specify the translations of typing, explicit
   computation witnesses and contradiction, with sufficient reflection or
   model soundness to support the intended consistency application. This
   must go beyond the all-true ground-Horn reduction: finite satisfiability
   alone does not exclude the translated contradiction.

# References

1. M. W. Bunder and W. J. M. Dekkers, *Are there Hilbert-style Pure Type Systems?*
   Logical Methods in Computer Science **4**(1:1), 2008, pp. 1--19.
   [doi:10.2168/LMCS-4(1:1)2008](https://doi.org/10.2168/LMCS-4(1:1)2008).
   [Open-access paper](https://lmcs.episciences.org/839/pdf).

2. R. B. Jones, *The Lambda-calculus, Combinatory Logic, and Type Systems*,
   1998--1999.
   [Conceptual overview](https://www.rbjones.com/rbjpub/logic/cl/cl017.htm#LSKICube).

3. A. Kaposi, A. Kovács and T. Altenkirch, *Constructing Quotient
   Inductive-Inductive Types*, Proceedings of the ACM on Programming
   Languages **3**(POPL), Article 2, 2019.
   [doi:10.1145/3290315](https://doi.org/10.1145/3290315).
   [Author manuscript](https://akaposi.github.io/finitaryqiit.pdf).

4. T. Altenkirch and A. Kaposi, *Type Theory in Type Theory using Quotient
   Inductive Types*, POPL 2016, pp. 18--29.
   [doi:10.1145/2837614.2837638](https://doi.org/10.1145/2837614.2837638).
   [Conference paper](https://akaposi.github.io/tt-in-tt.pdf);
   [expanded draft](https://akaposi.github.io/paper_ttintt.pdf), Sections 3.1--3.3.

5. L.-T. Chen, F. Nordvall Forsberg and T.-C. Tsai, *Can We Formalise Type
   Theory Intrinsically without Any Compromise? A Case Study in Cubical Agda*,
   CPP 2026. [doi:10.1145/3779031.3779090](https://doi.org/10.1145/3779031.3779090).
   [Author paper](https://fredriknf.com/papers/tt-qiirt_cpp2026.pdf);
   [Agda artifact](https://github.com/L-TChen/TTasQIIRT), inspected at
   commit `8db0830`. Sections 3--5 develop the natural-model QIIRT method.

6. M. Carneiro, *Metamath Zero: Designing a Theorem Prover Prover*, CICM 2020,
   LNCS **12236**, pp. 71--88.
   [doi:10.1007/978-3-030-53518-6\_5](https://doi.org/10.1007/978-3-030-53518-6_5).
   [MM0 specification](https://github.com/digama0/mm0/blob/0d414c0bfdaaeb7fea571895127abc1fa5a3d956/mm0.md).

7. S. Broda and L. Damas, *On combinatory complete sets of proper
   combinators*, Journal of Functional Programming **7**(6), 1997,
   pp. 593--612. [doi:10.1017/S0956796897002888](https://doi.org/10.1017/S0956796897002888).

8. [Original PTS-rk manuscript](../../latex/PTS-rk.tex), sections labelled
   `complexity`, `rep-bound`, `uniform-elementary` and `non-uniform-rep`, and the checked
   [normalization development](../NormalizationComplexity/README.md).
   See `BoundedNormalization`, `Runtime`, `EstimatorChecks` and `AutomaticBound`.

9. F. Pakhomov, *A weak set theory that proves its own consistency*,
   arXiv:1907.00877v2, 2019. [Paper](https://arxiv.org/pdf/1907.00877v2).
   Lemma 3.11 and Corollary 5.3 give the predicate-only transfer results;
   Appendix A relates the theory to Robinson's $R$.

10. B. van den Berg and M. den Besten, *Quadratic type checking for objective
    type theory*, arXiv:2102.00905, 2021.
    [Paper](https://arxiv.org/abs/2102.00905).
    Sections 2--3, especially Theorem 7.

%clearpage{}
%appendix{}

# Complete MM0 source {#app-mmzero}

The following listing is read from the same file as the excerpt in
Section <#sec-mmzero>. The highlighting follows MM0's declarations, comments and
dollar-delimited math strings.

%CombinatoryFullMMZeroListing{}
